% =====================================================================
%  paper-v7/sg1 --- SG1-1: the no-go theorem for the position--defect
%  decomposition  s_1 = F + G  with  s_1 := inf supp nu_w.
%
%  Self-contained: no \input, one file.  Build:
%      pdflatex -interaction=nonstopmode main.tex     (run twice)
%
%  English translation of the corpus note.  Every number, formula, label,
%  \ref and every [T]/[C]/[NG]/[E] chip is carried over unchanged.  The only
%  substantive edit is the grade of SG1-1, [T] -> [C] (see the convention
%  "conv:iii-strict"), together with the three limitations of the corpus and
%  the precise N09-1 relation, now written into the body.
%
%  House rules honoured here:
%   * labels are literal text  [T] / [C] / [NG] / [O]  (plus [E] for
%     machine-checked numbers); no \LT-style macros in this sub-paper;
%   * no completion stamp anywhere (CLAUDE.md sec.2)  ->  \qedsymbol empty;
%   * every assertion carries a source of the form file:line;
%   * Clay = 0%.
% =====================================================================
\documentclass[11pt]{article}

\usepackage[T1]{fontenc}
\usepackage[utf8]{inputenc}
\usepackage{lmodern}
\usepackage[margin=2.4cm,top=2.15cm,bottom=2.2cm]{geometry}
\usepackage{amsmath,amssymb,amsthm}
\usepackage{microtype}
\usepackage{array}
\usepackage[colorlinks=true,linkcolor=blue,urlcolor=blue]{hyperref}

% ---------- theorem environments, shared counter ----------
\newtheorem{theorem}{Theorem}
\newtheorem{lemma}[theorem]{Lemma}
\newtheorem{proposition}[theorem]{Proposition}
\newtheorem{corollary}[theorem]{Corollary}
\theoremstyle{definition}
\newtheorem{definition}[theorem]{Definition}
\newtheorem{convention}[theorem]{Convention}
\theoremstyle{remark}
\newtheorem{remark}[theorem]{Remark}

% ---------- no completion stamp (CLAUDE.md sec.2) ----------
\renewcommand{\qedsymbol}{}

% ---------- code / file:line citations ----------
\newcommand{\code}[1]{\texttt{\footnotesize\detokenize{#1}}}

% ---------- textual grade chips ----------
\newcommand{\grT}{\textsf{[T]}}
\newcommand{\grC}{\textsf{[C]}}
\newcommand{\grN}{\textsf{[NG]}}
\newcommand{\grO}{\textsf{[O]}}
\newcommand{\grE}{\textsf{[E]}}

\setlength{\parskip}{2pt}
\emergencystretch=2em

\begin{document}

\hypersetup{pdftitle={SG1-1: a no-go theorem for the position-defect decomposition},
  pdfsubject={covariant + lower side + finite-moment computable = empty},
  pdfkeywords={spectral measure, inf supp, N09-1, SSF/Krein-Birman, no-go}}
\title{\textbf{SG1-1}: a no-go theorem for the position--defect decomposition\\[3pt]
\large $\texorpdfstring{s_1=F+G}{s1=F+G}$, \
$\texorpdfstring{s_1:=\inf\operatorname{supp}\nu_w}{s1 := inf supp nu\_w}$, is not
reachable by a covariant, lower-bounded, finite-moment functional}
\author{Yang--Mills mass-gap corpus \quad \texttt{docs/paper-v7/sg1}}
\date{2026-11-25}
\maketitle

\vspace{-1.0em}
\begin{quote}\small
\textbf{Scope and grade.} This sub-paper states and proves one theorem,
\textbf{SG1-1} \grC{} --- the grade is \grC{}, not \grT: the proof is a
three-line argument whose Step~2 needs condition (iii) read as
\emph{continuity} in finitely many moments, and under the computability
reading of (iii) that step does not follow, so the step itself is \grC{}
(Convention~\ref{conv:iii-strict}). It localises what it excludes. It uses
only the measure-theoretic layer ($\nu\ge 0$, finite moments) and the
four-model slot table already fixed in the corpus. Labels are literal: \grT{}
= proved or identity, \grC{} = conditional, \grN{} = no-go for a named class,
\grO{} = exact missing step, \grE{} = machine-checked number
(\code{w7/sg1_G.log}, \texttt{py -3.12} + \texttt{mpmath 1.3.0}, exit $0$).
No completion stamp is attached to any proof. Clay $=0\%$. Only files under
\texttt{docs/paper-v7/sg1/} are written; no \code{git commit} is performed.
\end{quote}

% =====================================================================
\section[Statement]{Statement}\label{sec:statement}

\subsection[Carrier]{Carrier: $s_1$ is the bottom of the support of an $O$-channel measure}

\begin{definition}[spectral measure, weight, $s_1$]\label{def:s1}
Let $H\ge 0$ be self-adjoint with ground energy $E_0$, $\Omega$ its ground
vector, and $O=O^\ast$ a gauge-invariant operator. Put
\begin{equation}
\nu_w(B)\;=\;\langle\Omega|\,O\,P_H(B)\,O\,|\Omega\rangle
\;-\;|\langle\Omega|O|\Omega\rangle|^{2}\,\delta_0(B),
\qquad
\nu_w=\sum_{n} w_n\,\delta_{E_n-E_0},
\qquad
w_n=|\langle\Omega|O|n\rangle|^{2},
\label{eq:nuw}
\end{equation}
and
\begin{equation}
s_1(\nu_w)\;:=\;\inf\operatorname{supp}\nu_w .
\label{eq:s1}
\end{equation}
\end{definition}

\noindent
Sources: \code{ATOM2.md:27} for \eqref{eq:nuw}; \code{NS2.md:29} for
$w_n=|\langle\Omega|O|n\rangle|^{2}$; \code{NS2.md:31} for \eqref{eq:s1}.
Thus $s_1$ is the \emph{bottom of the support of the $O$-channel spectral
measure} --- a property of the pair $(\text{model},O)$, not of the model alone.

\begin{remark}[four limitations that travel with $s_1$]\label{rem:limits}
\mbox{}
\begin{enumerate}\itemsep1pt
\item \textbf{Per operator} \grT. $\nu$ is per operator: changing $O$ changes
the measure (\code{ATOM2.md:31}). On the same 2D Yang--Mills model,
$O=\tfrac14\operatorname{tr}F^2$ gives $\nu\equiv0$ (hence
$s_1=+\infty$, meaningless), while $O=\tfrac12\operatorname{tr}U$ gives an
atom of exact weight $1$ at $E_1=3g^2/8$ (\code{ATOM2.md:15,16,77--80}; quoted
at \code{SG1.md:36}). Every instance below therefore carries its $O$.
\item \textbf{$\tau_a$ is not a symmetry} \grT. There is no operator that
implements a shift of the spectrum; $\tau_a$ acts on $\nu$ only and not on the
data $(S_W,\beta,L,a)$ (\code{B05-DUAL.md:85}, quoted at \code{SG1.md:37}).
\item \textbf{$s_1=m$ is a physical premise} \grC, not a theorem:
\code{B05-DUAL.md:159} lists "position" and $m_1$ in the same cell and marks
the identification as \emph{equal to the target (lossless) $\Rightarrow$
circular}. It is used here only in that conditional reading.
\item \textbf{$s_1$ is not the model gap} \grT. On the model M6
(oscillator $\otimes$ free mode, $\omega=w=1$) one has $s_1=\omega=1$
(\code{ATOM2.md:44}) while the model spectral bottom is
$m_6=\min(\omega,2\pi^2/L^2)\to0$ (\code{N09-RESTATE.md:46,65}):
\begin{center}\footnotesize
\setlength{\tabcolsep}{4pt}%
\begin{tabular}{l|cccc}
$L$ & $8$ & $16$ & $32$ & $64$\\\hline
$m_6$ & $0.30842514$ & $0.077106284$ & $0.019276571$ & $0.0048191428$
\end{tabular}
\end{center}
\grE{} (\code{w7/sg1_G.log:92--95}); the three-point ladder $L=10,40,160$ gives
$m_6=0.1974$, $0.01234$, $0.000771\to0$ (\code{SG2.md:90}). Hence $s_1$ is
\emph{not} a function of the gap at the level
of values --- this is not a defect of the theorem below but its subject matter
(\S\ref{sec:what}).
\end{enumerate}
\end{remark}

\subsection{The three conditions}

Throughout, $\mathcal M$ denotes the class of non-negative Borel measures on
$[0,\infty)$ with finite moments $m_k(\nu):=\int\lambda^k\,d\nu(\lambda)$ for
$k\le N$ and with $s_1(\nu)=\inf\operatorname{supp}\nu\in\mathbb R$. (Spectral
measures \eqref{eq:nuw} form a subclass; the theorem is proved on all of
$\mathcal M$, hence \emph{a fortiori} on the physical subclass.) That
quantifier domain is fixed once and for all:

\begin{convention}[quantifier domain of (ii): the whole measure class]\label{conv:domain}
All three conditions below are quantified over the \emph{entire} class
$\mathcal M$ of non-negative Borel measures on $[0,\infty)$ with finite
moments up to the fixed $N$ and $s_1(\nu)\in\mathbb R$ --- i.e.\ over
\emph{all measure classes}, not merely over the spectral measures
\eqref{eq:nuw} that physics supplies. The statement is therefore
\emph{stronger} than its physical-class version: a functional that fails on
the larger class fails \emph{a fortiori} on the physical subclass, and the
emptiness proved below is an emptiness on all measure classes.
\end{convention}

\begin{definition}[G1 shift]\label{def:G1}
For $a\in\mathbb R$ set $(\tau_a\nu)(B):=\nu(B-a)$, i.e.\ $\tau_a\nu$ is the
energy-axis translate. For every $\nu\in\mathcal M$,
\begin{equation}
s_1(\tau_a\nu)\;=\;a+s_1(\nu).
\label{eq:s1cov}
\end{equation}
\end{definition}
\noindent\eqref{eq:s1cov} is \grT{} (\code{B05-DUAL.md:80}).

\begin{definition}[the three conditions]\label{def:three}
A functional $F:\mathcal M\to\mathbb R$ satisfies
\begin{enumerate}\itemsep2pt
\item[\textbf{(i)}] \textbf{G1-covariance}: $F(\tau_a\nu)=F(\nu)+a$ for all
$\nu\in\mathcal M$ and all $a\in\mathbb R$;
\item[\textbf{(ii)}] \textbf{lower side}: $F(\nu)\le s_1(\nu)$ for all
$\nu\in\mathcal M$;
\item[\textbf{(iii)}] \textbf{finite-moment computability}: there are
$N<\infty$ and a continuous $\Phi:\mathbb R^{N}\to\mathbb R$ with
$F(\nu)=\Phi\bigl(m_1(\nu),\dots,m_N(\nu)\bigr)$ for all $\nu\in\mathcal M$.
\end{enumerate}
\end{definition}

\begin{convention}[which shift: (i) is pinned to the G1 energy axis]\label{conv:G1}
Condition (i) must be read on the \emph{energy} axis (\textbf{G1}), not on the
space axis (\textbf{G2}): under a spatial translate $U(x)$ the object
$\inf\operatorname{supp}\nu$ is invariant, so a G2 reading makes (i)
\emph{empty} --- it constrains nothing --- and simultaneously puts it in direct
conflict with (ii) (\code{B05-DUAL.md:76--83}, quoted at \code{NS2.md:18}).
The conflict is an artefact of the reading, not of the mathematics. Hence (i)
is pinned to G1 here, and no G2 reading is used anywhere below.
\end{convention}

\begin{convention}[reading of (iii): fixed $N$, $\Phi$ continuous, dependence on $m_1,\dots,m_N$ only]\label{conv:iii-strict}
Condition (iii) is defined in exactly one way, and this is the only reading
used: there is an $N<\infty$, \emph{fixed independently of $\nu$}, and a
continuous $\Phi:\mathbb R^{N}\to\mathbb R$ such that
$F(\nu)=\Phi\bigl(m_1(\nu),\dots,m_N(\nu)\bigr)$ for every $\nu\in\mathcal M$.
If (iii) is instead read as \emph{computability} --- recursion-theoretic, or
"computable from $(S_W,\beta,L,a)$ by an algorithm" --- then Step~2 of the
proof of Theorem~\ref{thm:sg1}, namely the passage
$F\bigl(\nu_\varepsilon^{(a)}\bigr)\to F(\delta_a)$, which uses nothing but
continuity of $\Phi$, no longer follows, and the three-line proof collapses.
That step is therefore itself \grC{}, and consequently the true grade of
\textbf{SG1-1} is \grC{}, not \grT. This is the only reason for the
downgrade, and it is stated here rather than buried in a remark.
\end{convention}

\begin{remark}[on (iii)]\label{rem:iii}
Condition (iii) is \emph{finite-moment continuity}, not computability in the
recursion-theoretic sense: it says $F$ depends on $\nu$ only through finitely
many moments $(m_1,\dots,m_N)$ with $N$ fixed independently of $\nu$, and
continuously. Every functional obtained by inserting finitely many moments into
a closed expression is of this form. The distinction is not cosmetic: the
proof of Theorem~\ref{thm:sg1} uses exactly continuity, and nothing else ---
so a "computable but discontinuous in the moment topology" $F$ is outside the
class killed below, and that is the one gap the theorem leaves
(Remark~\ref{rem:sharp-iii} and \S\ref{sec:what}).
\end{remark}

\subsection{The theorem}

\begin{theorem}[SG1-1 \grC]\label{thm:sg1}
\mbox{}
\[
\textbf{(i)}\ \text{G1-covariant}\quad\wedge\quad
\textbf{(ii)}\ F\le s_1\quad\wedge\quad
\textbf{(iii)}\ \text{finite-moment computable}
\quad=\quad\varnothing .
\]
Equivalently: no functional $F:\mathcal M\to\mathbb R$ satisfies
\ref{def:three}(i), (ii) and (iii) simultaneously, where (i) is the G1 energy
shift (Convention~\ref{conv:G1}), (ii) quantifies over the whole class
$\mathcal M$ (Convention~\ref{conv:domain}) and (iii) is the strict reading of
Convention~\ref{conv:iii-strict}. The grade is \grC, because the last step of
the proof is conditional on that reading.
\end{theorem}

\noindent
The corpus statement, transcribed from \code{SG1.md:62--67}
(markdown emphasis and backticks are markup; mathematical symbols are restored
in math mode; no word is changed, all Chinese wording is translated):

\begin{quote}\small
\textbf{SG1-1 \grC}: \textbf{(i)} G1-covariance $\wedge$ \textbf{(ii)}
$F\le s_1$ $\wedge$ \textbf{(iii)} finite-moment computable $=\varnothing$\\
Proof: the atom family $\nu_\varepsilon^{(a)}$ makes the moments drift by
$O(\varepsilon)$ while $s_1:1\to\varepsilon$\\
\hspace*{2.1em}machine check $a=0.5/1/2/5$ all \textsf{NO} (\code{sg1_G.log})\\
$\Rightarrow$ the only escape left is the \textbf{full-moment limit}
$\Rightarrow$ the only reachable object is $F=s_1$
$\Rightarrow$ $G\equiv0$ $\Rightarrow$ collides with N09-1\\
and $G=s_1$ is exactly the \textbf{SSF/Krein--Birman
$\operatorname{tr}\bigl(f(H)-f(H_0)\bigr)$ body} (\code{NS2:116})
\end{quote}

\noindent
\textbf{One transcription departure, declared.} The corpus text prints
\grT{} for SG1-1 at \code{SG1.md:62}; it is printed \grC{} here, in accordance
with Convention~\ref{conv:iii-strict}. This is the only place where the
transcription departs from the source, and the departure is the grade chip
only: every symbol, number and connective is preserved.

% =====================================================================
\section[Proof]{Proof}\label{sec:proof}

\subsection{The atom family}

\begin{lemma}[thin atom: all moments converge, $s_1$ jumps]\label{lem:atom}
Fix $a\ge0$, $\lambda>0$ and $0<\varepsilon<\min(a,2\lambda)$. Put
\begin{equation}
\nu_\varepsilon^{(a)}\;=\;\frac{\varepsilon}{2\lambda}\,\delta_\varepsilon
\;+\;\Bigl(1-\frac{\varepsilon}{2\lambda}\Bigr)\,\delta_a
\;\in\;\mathcal M .
\label{eq:atomfam}
\end{equation}
Then
\begin{equation}
s_1\bigl(\nu_\varepsilon^{(a)}\bigr)=\varepsilon,
\qquad
m_k\bigl(\nu_\varepsilon^{(a)}\bigr)
=\frac{\varepsilon}{2\lambda}\,\varepsilon^{k}
+\Bigl(1-\frac{\varepsilon}{2\lambda}\Bigr)a^{k}
\;\xrightarrow[\ \varepsilon\downarrow0\ ]{}\; a^{k}=m_k(\delta_a)
\quad(k\ge1).
\label{eq:atommom}
\end{equation}
\end{lemma}

\noindent
The family is the same thin-atom shape used against H1
(\code{A3-VERDICT.md:61--69}); here the weight $\varepsilon/2\lambda$ is
bookkeeping only --- the proof below needs $w_\varepsilon\to0$, nothing more.

\subsection[Proof of SG1-1]{Proof of Theorem~\ref{thm:sg1}}

\begin{proof}
Let $F$ satisfy \ref{def:three}(i)--(iii), with data $N,\Phi$. Fix $a\ge0$ and
let $\varepsilon\downarrow0$ along \eqref{eq:atomfam}. Three steps.

\smallskip
\noindent\emph{Step 1 (moments converge).} By \eqref{eq:atommom}, the moment
vector $\bigl(m_1,\dots,m_N\bigr)\bigl(\nu_\varepsilon^{(a)}\bigr)$ tends to
that of the single atom $\delta_a\in\mathcal M$.

\smallskip
\noindent\emph{Step 2 (lower side kills $F$ on atoms) \grC.} Continuity of
$\Phi$ in (iii) together with Step 1 gives
$F\bigl(\nu_\varepsilon^{(a)}\bigr)\to F(\delta_a)$, while (ii) gives
$F\bigl(\nu_\varepsilon^{(a)}\bigr)\le s_1\bigl(\nu_\varepsilon^{(a)}\bigr)=\varepsilon$.
Hence
\begin{equation}
F(\delta_a)\;\le\;0\qquad\text{for every }a\ge0 .
\label{eq:step2}
\end{equation}
The chip \grC{} sits on this step alone: it is exactly the step that a
computability reading of (iii) voids (Convention~\ref{conv:iii-strict}).

\smallskip
\noindent\emph{Step 3 (covariance kills atoms).} $(\tau_a\delta_0)(B)=\delta_0(B-a)=\delta_a(B)$,
so by (i) applied at $\nu=\delta_0$,
\begin{equation}
F(\delta_a)=F(\tau_a\delta_0)=F(\delta_0)+a \qquad(a\ge0).
\label{eq:step3}
\end{equation}
Combining \eqref{eq:step2} and \eqref{eq:step3},
\[
F(\delta_0)+a\;\le\;0\qquad\text{for every }a\ge0,
\]
which is false for any $a>|F(\delta_0)|$. Hence no such $F$ exists and the
intersection is empty (conditionally on the reading of (iii) fixed in
Convention~\ref{conv:iii-strict}).
\end{proof}

\begin{remark}[what the proof does and does not use]\label{rem:cheap}
\mbox{}
\begin{itemize}\itemsep1pt
\item It uses \emph{one} moment ($N\ge1$) and continuity, no more. Hence it
already kills every $F$ of the form $\Phi(m_1)$ with $\Phi$ continuous --- in
particular every "position estimate built out of the first moment".
\item It uses neither positivity of $F$ nor compact support nor gauge input:
(ii) is a two-sided squeeze against a quantity that jumps, and a single jump
suffices.
\item Conversely it does \emph{not} touch functionals that are discontinuous
in the moment topology --- and $s_1$ itself is exactly such a functional
(\eqref{eq:atommom} is a sequence with all moments convergent and
$s_1$ jumping from $\varepsilon$ to $a$). This is the precise content of the
dichotomy in Corollary~\ref{cor:dich}.
\end{itemize}
\end{remark}

\subsection[Machine check]{Machine check (\grE, \code{w7/sg1_G.log})}

Script: \code{research/v3/species/w7/sg1_G.py}; log:
\code{research/v3/species/w7/sg1_G.log}; status: \texttt{py -3.12},
\texttt{mpmath 1.3.0}, exit $0$ (\code{SG1.md:158}). Two tables are relevant.

\smallskip
\noindent\emph{(a) Moments drift by $O(\varepsilon)$ while $s_1$ moves by
$O(1)$} (base $\nu=\delta_{s_2}$, $s_2=1$, $\lambda=1$;
\code{sg1_G.log:37--42}):

\begin{center}\footnotesize
\begin{tabular}{r|r|r|r|r|r}
$\varepsilon$ & $s_1(\nu_\varepsilon)$ & $w_\varepsilon$ & $|\Delta m_1|$ &
$|\Delta m_2|$ & $|\Delta m_3|$\\\hline
$10^{-1}$ & $0.1$ & $0.05$ & $0.045$ & $0.0495$ & $0.04995$\\
$10^{-2}$ & $0.01$ & $0.005$ & $0.00495$ & $0.0049995$ & $0.00499999$\\
$10^{-6}$ & $10^{-6}$ & $5\cdot10^{-7}$ & $4.99999\cdot10^{-7}$ &
$5\cdot10^{-7}$ & $5\cdot10^{-7}$
\end{tabular}
\end{center}

\smallskip
\noindent
So $s_1$ drops $1\to\varepsilon$ (an $O(1)$ change) while every moment drifts
by $O(\varepsilon)$: finitely many moments cannot see the jump.

\smallskip
\noindent\emph{(b) The four-row inconsistency at $\varepsilon=10^{-6}$}
(\code{sg1_G.log:55--60}): for each $a$, (ii) forces $F\le10^{-6}$ while (i)
forces $F=a$.

\begin{center}\footnotesize
\begin{tabular}{r|r|r|r|r|c}
$a$ & $s_1(\nu_\varepsilon^{(a)})$ & $|\Delta m_1|$ vs $\delta_a$ &
(ii) forces $F\le$ & (i) forces $F=$ & consistent?\\\hline
$0.5$ & $1.0\cdot10^{-6}$ & $2.49999\cdot10^{-7}$ & $1.0\cdot10^{-6}$ & $0.5$ & \textsf{NO}\\
$1.0$ & $1.0\cdot10^{-6}$ & $4.99999\cdot10^{-7}$ & $1.0\cdot10^{-6}$ & $1.0$ & \textsf{NO}\\
$2.0$ & $1.0\cdot10^{-6}$ & $9.99999\cdot10^{-7}$ & $1.0\cdot10^{-6}$ & $2.0$ & \textsf{NO}\\
$5.0$ & $1.0\cdot10^{-6}$ & $2.5\cdot10^{-6}$ & $1.0\cdot10^{-6}$ & $5.0$ & \textsf{NO}
\end{tabular}
\end{center}

\noindent
All four rows are \textsf{NO}, as Theorem~\ref{thm:sg1} requires. Truncating
the moment list is precisely what the atom family kills: a finer grid of $a$
does not help, because the obstruction is the jump, not the resolution.

\subsection{The dichotomy, and where the survivor lands}

\begin{corollary}[dichotomy]\label{cor:dich}
The only way past Theorem~\ref{thm:sg1} is to abandon (iii): one must use the
\emph{full} moment sequence, i.e.\ the measure itself. On compact support the
full moment sequence determines $\nu$ uniquely, so the object reached this way
carries no less information than $\nu$; the only $(i)+(ii)$ object the corpus
can name in that region is $s_1$ itself, whence
\begin{equation}
F=s_1\qquad\Longrightarrow\qquad G:=s_1-F\equiv0 .
\label{eq:dich}
\end{equation}
\end{corollary}

\begin{remark}[window of the corollary --- it is a dichotomy, not an
impossibility theorem]\label{rem:sharp-iii}
\code{SG1.md:143} states the window verbatim: the conclusion of
Theorem~\ref{thm:sg1} covers only $F$ continuous in finitely many moments; the
full-moment route can be taken, but then "without consulting the spectrum"
fails. ``This is a dichotomy, not an absolute impossibility'' \grT. The
mechanism is the one isolated in Remark~\ref{rem:cheap}: $s_1$ is
discontinuous in the moment topology, and the escape region is the class of
functionals discontinuous at the atom limit. The corollary's identification of
the extremal survivor with $s_1$ is an identification of \emph{reachability},
and it is the weakest link of this section --- it is not a classification
theorem for that class.
\end{remark}

\begin{remark}[signature of the survivor: it lands inside N09-1's
domain]\label{rem:sig}
$G\equiv0$ has the four-model signature $(0,0,0,\text{undef})$ --- free field
$0$, 2D Yang--Mills $0$, M6 $0$, single plaquette undefined --- hence it
\emph{passes} the slot-value test $4/4$ and therefore lies in
$\operatorname{dom}(\mathrm{N09\text{-}1})=\{X:X\ \text{passes }\text{\textcircled{1}}\}$
(\code{Q03-TARGET.md:87}; slot table \code{N09-RESTATE.md:49},
\code{N12-VERDICT.md:36}). N09-1 then returns $G=f(m)$. So the survivor
produces no object beyond the target: the "defect" is the zero functional
(degenerate end $f\equiv0$), and $F=s_1$ is the position itself. See
\S\ref{sec:what} for the precise form of this "collision".
\end{remark}

% =====================================================================
\section[Sharpness]{Sharpness: dropping any one of the three conditions leaves solutions}\label{sec:sharp}

\begin{proposition}[mutual irredundancy]\label{prop:sharp}
Each of the three conditions is necessary: dropping any single one leaves a
non-empty class of solutions, and each witness below fails exactly one
condition and satisfies the other two.
\end{proposition}

\begin{center}\footnotesize
\begin{tabular}{@{}l|l|c|c|c|l@{}}
drop & witness $F$ & (i) & (ii) & (iii) & mechanism / source\\\hline
(i) & $F\equiv0$ & $\times$ & $\checkmark$ & $\checkmark$ &
$F(\tau_a\nu)-F(\nu)=0\ne a$; $0\le s_1$; $\Phi\equiv0$\\
(i) & $\mathrm{Var}$, $f_1^2$, $m_2$ & $\times$ & model-dep. & $\checkmark$ &
shifts $0,0,0,0$ at $a=0.1/0.5/1/2$ (\code{sg1_G.log:17--20})\\
(ii) & $F=\overline\lambda:=m_1/m_0$ & $\checkmark$ & $\times$ & $\checkmark$ &
$\overline\lambda-s_1\ge0$, equality iff $\nu$ is one atom\\
(iii) & $F=s_1$ & $\checkmark$ & $\checkmark$ (equality) & $\times$ &
$s_1(\tau_a\nu)=a+s_1(\nu)$; atom family \eqref{eq:atommom}
\end{tabular}
\end{center}

\noindent
Details, one line each.

\smallskip
\noindent\textbf{Dropping (i) $\Rightarrow$ shape.} $F\equiv0$ is covariant-free,
finite-moment ($N=0$) and below $s_1$ on $\mathcal M\subset$ measures on
$[0,\infty)$. The price is that $G:=s_1-F$ is then \emph{not} G1-invariant:
$G(\tau_a\nu)=G(\nu)+a$, so the decomposition of \code{NS2-0} collapses --- the
very identity $F=s_1-G$ with $G$ a \emph{shape} functional requires (i). The
same verdict covers the shape candidates: $\mathrm{Var}$, $f_1^2$ and $m_2$
all have zero shift under $\tau_a$ (\code{sg1_G.log:17--20}), hence satisfy
(i) for no $a\ne0$; the census of covariant functionals among the tested
candidates is exactly $\{s_1,\ \overline\lambda\}$ (\code{sg1_G.log:22}).

\smallskip
\noindent\textbf{Dropping (ii) $\Rightarrow$ the mean.}
$F=\overline\lambda=m_1/m_0$ is G1-covariant ($+a$ exactly, at
$a=0.1/0.5/1/2$: \code{sg1_G.log:16}) and finite-moment ($N=2$, $\Phi$ smooth
on $m_0>0$). It fails (ii) on the class: $\overline\lambda\ge s_1$ with
equality iff $\nu$ is a single atom, so on
$\nu=\tfrac12\delta_0+\tfrac12\delta_1$ one has
$\overline\lambda-s_1=0.5-0=0.5>0$ (\code{sg1_G.log:12,30}). This is the corpus's
NG-3 (\code{B05-DUAL.md:141--143}): the mean is on the \emph{wrong side}, so it
can only give upper bounds --- the opposite of what a lower bound on the
position needs.

\smallskip
\noindent\textbf{Dropping (iii) $\Rightarrow$ $s_1$ itself.}
$F=s_1$ is perfectly covariant (\eqref{eq:s1cov}) and satisfies (ii) with
equality. It is not finite-moment computable: \eqref{eq:atommom} exhibits two
measures whose first $N$ moments agree to $O(\varepsilon)$ while their $s_1$
differ by $1-\varepsilon$. So the only exact witness of (i)${}+{}$(ii) is the
target, and the corpus's own summary of this degeneration is
\code{SG1.md:105}: "either it degenerates into the position itself, or one has
to hand over $G$" (translated from the corpus wording).

\smallskip
\noindent
Hence the emptiness in Theorem~\ref{thm:sg1} is a genuine three-way
incompatibility, not one condition carrying the whole theorem, and the
sharpness table also displays the cost of each concession: shape (drop (i)) is
available but carries no certificate direction; the mean (drop (ii)) is
available but on the upper side; the position (drop (iii)) is available but is
the target.

% =====================================================================
\section[What it excludes]{What it excludes}\label{sec:what}

\subsection[The sixth species]{The sixth species: the specification \code{NS2-S} cannot be met}

The corpus's sixth-species specification (\code{NS2.md:59--69}) asks for a pair
of computable functionals $(F,G)$ with
\begin{quote}\footnotesize
(1) $s_1=F+G$ (exact identity), \quad
(2) $F=F(S_W,\beta,L,a)$ computable \emph{without} the spectrum, \quad
(3) $F(\tau_a\nu)=F(\nu)+a$, $G(\tau_a\nu)=G(\nu)$, \quad
(4) $G\ge0$ computable and not a renaming of $H1/\mathrm{Var}/G''/f_1^2$, \quad
(5) $F(\mathrm{M6})=0$, $G(\mathrm{M6})=\omega\ne0=m_6$, \quad
(6) $F\ge c_0$ $(L,a)$-uniformly $\Rightarrow$ $s_1\ge c_0$.
\end{quote}
Items (1)--(3) are already fixed as \grT{} algebra (\code{NS2-0},
\code{NS2.md:38--47}): given any G1-invariant $G\ge0$, $F:=s_1-G$ satisfies
both premises automatically, so the identity carries no information and the
only real question is whether an admissible $F$ exists at all
(\code{SG1.md:46}). Theorem~\ref{thm:sg1} answers it in the negative for the
finite-moment branch:

\begin{itemize}\itemsep2pt
\item $(2)\wedge(3)\wedge[\,F\ \text{finite-moment, i.e.\ in (iii)}\,]=\varnothing$
(Theorem~\ref{thm:sg1}, with $c_0>0$ in (6) playing no role --- the obstruction
is already at $F\le s_1$);
\item the two ways out are exactly (6)'s degenerations: either $F=s_1$
($G\equiv0$, no defect, and (6) becomes the target statement itself), or $F$
violates (i) or (ii) and then \code{NS2-0}'s split of the certificate
direction is void.
\end{itemize}

\noindent
One conditional step must be flagged, because it is not part of
Theorem~\ref{thm:sg1}: the passage from "(2) $F$ is computed from
$(S_W,\beta,L,a)$ without the spectrum" to "(iii) $F$ is continuous in finitely
many moments" is a \emph{reading}, marked \grC{} here. \code{SG1.md:132}
asserts it flatly ("any $F$ computed directly from $S_W$ is a function of
finitely many moments", translated); this paper keeps the assertion but labels
it \grC, since a computable functional could in principle be discontinuous in
the moment topology --- which is exactly the gap in
Remark~\ref{rem:sharp-iii} and the reason for the grade \grC{} carried by
Theorem~\ref{thm:sg1} itself.

\subsection[The specification]{The specification $F=s_1-G$: the decomposition is a tautology plus
a hole}

Assembling the above: the decomposition $s_1=F+G$ is \grT{} but content-free,
and Theorem~\ref{thm:sg1} says the pair cannot be made non-vacuous on the
finite-moment side. In other words, for the specification to succeed one needs
an $F$ that is simultaneously
\begin{enumerate}\itemsep1pt
\item[(a)] G1-covariant, (b) below the position, (c) computable without the
spectrum, (d) positive uniformly;
\end{enumerate}
and (a)${}\wedge{}$(b)${}\wedge{}$(c) already has no solution. The only survivor
of (a)${}\wedge{}$(b) is $F=s_1$, for which (c) fails and (d) is the target.
So the residual content of the specification is: \emph{prove the position is
positive} --- verbatim the Clay-side statement, not a new object.

\subsection[SG2]{The escape of \S99.3 is a false positive, and so is $s_1$}

The corpus's \S99.3 promoted $F=s_1-G$ to "the first candidate in 60+ rounds to
escape N09-1", on the ground that $G(\mathrm M6)=\omega\ne0$ while $m_6\to0$.
\code{SG2.md} kills that reading in two steps, both inherited here:

\smallskip
\noindent\emph{(1) $G$ violates the M6 slot value.} The slot-value test demands
$X(\mathrm M6)=0$ \emph{exactly} (\code{N09-RESTATE.md:49}: "must take exactly
0", translated); $G(\mathrm M6)=\omega=1\ne0$ (\code{ATOM2.md:44}). This is a
false positive of exactly the $\Delta_0$ type, whose four-model row reads
"M6 positive (false positive) $\times$" (\code{N12-VERDICT.md:13}). Same type,
same verdict.

\smallskip
\noindent\emph{(2) $s_1$ itself violates it.} With the specification's own
choice $F(\mathrm M6)=0$ (\code{NS2.md:65}),
\[
s_1(\mathrm M6)=F(\mathrm M6)+G(\mathrm M6)=0+\omega=1\;>\;0=m_6 ,
\]
with $\omega=1$ (\code{ATOM2.md:44}) and $m_6\to0$ under $L\to\infty$
(\code{N09-RESTATE.md:65});
so the object the specification wants to bound from below is \emph{itself} a
false positive on the M6 slot. Splitting $s_1$ into $F+G$ does not move the
defect off the object: whichever summand carries $\omega$, the sum carries it
(\code{SG2.md:99--104}). Hence
\begin{quote}\small
\textbf{escape $\equiv$ expulsion} (\code{SG2.md:135--138}; translated):
escaping $\operatorname{dom}(\mathrm{N09\text{-}1})$ and failing the four-model
test are the same statement, because the domain \emph{is} the set of
test-passers (\code{Q03-TARGET.md:87}).
\end{quote}

\subsection{Precise relation to N09-1, and what remains open}

\begin{remark}[SG1-1 and N09-1: no implication in either direction; they meet only at a degenerate end]\label{rem:n09rel}
\mbox{}
\begin{itemize}\itemsep1pt
\item \textbf{No implication either way.} N09-1
(\code{N09-RESTATE.md:44--53}, \grT{} conditional on the M6 specification and
the $L\to\infty$ reading) quantifies over \emph{quantities on models} that pass
the slot-value test; SG1-1 quantifies over \emph{functionals on measures}
(Convention~\ref{conv:domain}). Neither implies the other.
\item \textbf{Where they meet.} Only at the survivor: $G\equiv0$ passes $4/4$
and lands inside $\operatorname{dom}(\mathrm{N09\text{-}1})$
(Remark~\ref{rem:sig}). That contact is a \emph{degenerate end}
($F=s_1\Rightarrow G\equiv0$), not an implication: it is one shared boundary
point of two different statements, not one statement entailing the other.
\item \textbf{The word "collision" is not an implication.} In the corpus's
phrasing the branch $F=s_1$ "collides with N09-1". Precisely: N09-1's
conclusion is $G=f(m)$ with $f(0)=0$ and $f>0$; the survivor is $f\equiv0$, so
the collision is at the \emph{degenerate end} of that family. What is excluded
is the existence of a \emph{non-trivial} defect $G$; the zero defect is not an
object at all. This refinement should be carried whenever the corpus's
collision phrase is quoted.
\item \textbf{Grade bookkeeping.} Since SG1-1 is \grC{} and N09-1 is \grT{}
(conditional on its own premises), no transfer of grade in either direction is
available, which is a second, independent reason why no implication holds.
\end{itemize}
\end{remark}

\smallskip
\noindent\textbf{What remains open} \grO. The class not touched by
Theorem~\ref{thm:sg1} is the full-moment / moment-discontinuous class. There
the corpus's named object is the SSF/Krein--Birman relative trace
$\operatorname{tr}\bigl(f(H)-f(H_0)\bigr)=\int f'\xi$ with
$\xi_{H+a}(\lambda)=\xi_H(\lambda-a)$, which is G1-covariant and whose bottom is
transported losslessly --- the corpus's own verdict is that this \emph{is} the
circular body of (i) and simultaneously a transport argument, i.e.\ the class
with $0/19$ successes (\code{NS2.md:116}, \code{NS2.md:83}). So the open region
is exactly "the position, transported", and it is where the Clay-side work
already lives. No new object is produced there by this paper, and none is
claimed.

\medskip
\noindent\textbf{Honest summary.} This sub-paper adds no new quantity and no new
numerical evidence: \eqref{eq:atomfam}--\eqref{eq:step3} is a three-line
argument, the tables are re-readings of \code{sg1_G.log}, and the fourth
section is bookkeeping over already-ruled corpus items. What is added is the
isolation of SG1-1 as a statement with its three conditions made explicit ---
including the pinning of (i) to the G1 energy axis, the quantifier domain of
(ii) as the whole measure class, and the strict reading of (iii) that fixes the
grade at \grC{} --- its sharpness witnesses, the exact N09-1 relation (degenerate
contact, no implication), and the exact boundary of what it does \emph{not}
exclude. Clay $=0\%$; no completion stamp.

\end{document}
